\documentclass[10pt,a4paper]{amsart}
\usepackage{amsmath,amssymb,amsthm,graphicx}
\usepackage{verbatim}
\newenvironment{reference}{\comment}{\endcomment}
%\newenvironment{reference}{}{}
\newenvironment{slogan}{\comment}{\endcomment}
\newenvironment{history}{\comment}{\endcomment}

% For commutative diagrams we use Xy-pic
\usepackage[all]{xy}

% We use 2cell for 2-commutative diagrams.
\xyoption{2cell}
% Theorem environments.
%
\theoremstyle{plain}
\newtheorem{theorem}[subsection]{Theorem}
\newtheorem{proposition}[subsection]{Proposition}
\newtheorem{lemma}[subsection]{Lemma}

\theoremstyle{definition}
\newtheorem{definition}[subsection]{Definition}
\newtheorem{example}[subsection]{Example}
\newtheorem{exercise}[subsection]{Exercise}
\newtheorem{situation}[subsection]{Situation}

\theoremstyle{remark}
\newtheorem{remark}[subsection]{Remark}
\newtheorem{remarks}[subsection]{Remarks}

\numberwithin{equation}{subsection}

% Macros
%
\def\lim{\mathop{\mathrm{lim}}\nolimits}
\def\colim{\mathop{\mathrm{colim}}\nolimits}
\def\Spec{\mathop{\mathrm{Spec}}}
\def\Hom{\mathop{\mathrm{Hom}}\nolimits}
\def\Ext{\mathop{\mathrm{Ext}}\nolimits}
\def\SheafHom{\mathop{\mathcal{H}\!\mathit{om}}\nolimits}
\def\SheafExt{\mathop{\mathcal{E}\!\mathit{xt}}\nolimits}
\def\Sch{\mathit{Sch}}
\def\Mor{\mathop{\mathrm{Mor}}\nolimits}
\def\Ob{\mathop{\mathrm{Ob}}\nolimits}
\def\Sh{\mathop{\mathit{Sh}}\nolimits}
\def\NL{\mathop{N\!L}\nolimits}
\def\CH{\mathop{\mathrm{CH}}\nolimits}
\def\proetale{{pro\text{-}\acute{e}tale}}
\def\etale{{\acute{e}tale}}
\def\QCoh{\mathit{QCoh}}
\def\Ker{\mathop{\mathrm{Ker}}}
\def\Im{\mathop{\mathrm{Im}}}
\def\Coker{\mathop{\mathrm{Coker}}}
\def\Coim{\mathop{\mathrm{Coim}}}

% Boxtimes
%
\DeclareMathSymbol{\boxtimes}{\mathbin}{AMSa}{"02}

%
% Macros for moduli stacks/spaces
%
\def\QCohstack{\mathcal{QC}\!\mathit{oh}}
\def\Cohstack{\mathcal{C}\!\mathit{oh}}
\def\Spacesstack{\mathcal{S}\!\mathit{paces}}
\def\Quotfunctor{\mathrm{Quot}}
\def\Hilbfunctor{\mathrm{Hilb}}
\def\Curvesstack{\mathcal{C}\!\mathit{urves}}
\def\Polarizedstack{\mathcal{P}\!\mathit{olarized}}
\def\Complexesstack{\mathcal{C}\!\mathit{omplexes}}
% \Pic is the operator that assigns to X its picard group, usage \Pic(X)
% \Picardstack_{X/B} denotes the Picard stack of X over B
% \Picardfunctor_{X/B} denotes the Picard functor of X over B
\def\Pic{\mathop{\mathrm{Pic}}\nolimits}
\def\Picardstack{\mathcal{P}\!\mathit{ic}}
\def\Picardfunctor{\mathrm{Pic}}
\def\Deformationcategory{\mathcal{D}\!\mathit{ef}}
\usepackage{cleveref}
\crefname{originallemma}{original lemma}{original lemmas}
\Crefname{originallemma}{Original lemma}{Original lemmas}
\makeatletter
\newlabel{lemma-relative-dualizing-complex-algebra}{{0E2U}{source}{}{sourcepointer-1}{}}
\newlabel{lemma-relative-dualizing-complex-algebra@cref}{{[originallemma][0][]0E2U}{[1][0][]source}{}{}{}}
\makeatother
\crefname{originallemma}{original lemma}{original lemmas}
\Crefname{originallemma}{Original lemma}{Original lemmas}
\makeatletter
\newlabel{dualizing-lemma-uniqueness-relative-dualizing}{{0E2D}{source}{}{sourcepointer-2}{}}
\newlabel{dualizing-lemma-uniqueness-relative-dualizing@cref}{{[originallemma][0][]0E2D}{[1][0][]source}{}{}{}}
\makeatother
\crefname{originallemma}{original lemma}{original lemmas}
\Crefname{originallemma}{Original lemma}{Original lemmas}
\makeatletter
\newlabel{lemma-relative-dualizing-RHom}{{0E2V}{source}{}{sourcepointer-3}{}}
\newlabel{lemma-relative-dualizing-RHom@cref}{{[originallemma][0][]0E2V}{[1][0][]source}{}{}{}}
\makeatother
\crefname{originallemma}{original lemma}{original lemmas}
\Crefname{originallemma}{Original lemma}{Original lemmas}
\makeatletter
\newlabel{cohomology-lemma-vanishing-and-glueing}{{0D66}{source}{}{sourcepointer-4}{}}
\newlabel{cohomology-lemma-vanishing-and-glueing@cref}{{[originallemma][0][]0D66}{[1][0][]source}{}{}{}}
\makeatother
\title{A relative dualizing complex is unique up to a unique isomorphism}
\author{The Stacks Project Authors}
\date{stacks-project:duality.tex:local-snapshot-20261012; GNU FDL 1.2 or later}
\begin{document}
\maketitle
\noindent\textit{Editorial note.} All delivered mathematical text is the author's own, copied from the pinned Stacks Project chapter file duality.tex; no statement or proof step was rewritten, reordered or removed.\par\medskip
\noindent\textit{Editorial note.} Preamble provenance: the pinned chapter file opens with an input of the shared preamble and carries no preamble text of its own, so the project's single Stacks-wide preamble staged at sources/stacks-duality/preamble.tex is excerpted verbatim under the recipe fixed by LANE\_BRIEF 22.1, as three spans, lines 8-13 (verbatim package and the reference/slogan/history environments), lines 14-18 (Xy-pic setup) and lines 154-217 (the native theorem declarations and mathematical macros). Nothing was added, removed or redefined; the wrapper class and the amsmath/amssymb/amsthm/graphicx base packages are supplied by the assembler.\par\medskip
\noindent\textit{Editorial note.} The counted statement and proof contain no cross reference, no citation and no included file, so the delivered body is self contained.\par\medskip
\noindent\textit{Editorial note.} Canonical location of the counted proposition: Stacks tag 0E2W (sources/stacks-duality/pinned-original-source.tex line 7334).\par\medskip
\noindent\textit{Allowed standard prerequisites (editorial).}\begin{itemize}
\item The affine algebraic uniqueness of the relative dualizing complex. Same chapter, Lemma labelled lemma-relative-dualizing-complex-algebra, cited at proof lines 7341-7344, and Dualizing Complexes, Lemma labelled dualizing-lemma-uniqueness-relative-dualizing, cited at proof lines 7344-7346. Provides the isomorphisms on affine opens.
\item The automorphism group and negative Ext vanishing for a relative dualizing complex. Same chapter, Lemma labelled lemma-relative-dualizing-RHom, cited at proof lines 7358-7360 and 7366-7368. Shows there is at most one normalised isomorphism and gives the vanishing needed to glue.
\item Glueing of morphisms from the vanishing of negative Ext groups. Cohomology, Lemma labelled cohomology-lemma-vanishing-and-glueing, cited at proof lines 7368-7372. Produces the global morphism agreeing with the local ones.
\item The generator property of the comparison map defining a relative dualizing complex. The invertibility of the Delta\_*O\_X-module generated by eta and alpha(xi), used at proof lines 7350-7356. Allows alpha to be normalised by an invertible section.
\end{itemize}\medskip
\ifdefined\hypertarget\hypertarget{sourcepointer-1}{}\fi
\noindent\textit{Source reference.} Original lemma 0E2U: sources/stacks-duality/pinned-original-source.tex, source lines 7271--7271. Same-chapter pointer into the author environment lemma at source line 7271; invoked at proof line 7344 of the counted proof. It is retained only as a navigation pointer; every load-bearing step is inside the delivered span.\par\medskip
\ifdefined\hypertarget\hypertarget{sourcepointer-2}{}\fi
\noindent\textit{Source reference.} Original lemma 0E2D: sources/stacks-chow/tags.txt, source lines 17603--17603. Cross-chapter standard prerequisite located in the project tag list staged in the local snapshot; invoked at proof line 7346 of the counted proof. It is retained only as a navigation pointer; every load-bearing step is inside the delivered span.\par\medskip
\ifdefined\hypertarget\hypertarget{sourcepointer-3}{}\fi
\noindent\textit{Source reference.} Original lemma 0E2V: sources/stacks-duality/pinned-original-source.tex, source lines 7318--7318. Same-chapter pointer into the author environment lemma at source line 7318; invoked at proof line 7360 of the counted proof. It is retained only as a navigation pointer; every load-bearing step is inside the delivered span.\par\medskip
\ifdefined\hypertarget\hypertarget{sourcepointer-4}{}\fi
\noindent\textit{Source reference.} Original lemma 0D66: sources/stacks-chow/tags.txt, source lines 16511--16511. Cross-chapter standard prerequisite located in the project tag list staged in the local snapshot; invoked at proof line 7371 of the counted proof. It is retained only as a navigation pointer; every load-bearing step is inside the delivered span.\par\medskip
\begin{lemma}
\label{lemma-uniqueness-relative-dualizing}
Let $X \to S$ be a morphism of schemes which is flat and
locally of finite presentation. If $(K, \xi)$ and $(L, \eta)$
are two relative dualizing complexes on $X/S$, then there is a unique
isomorphism $K \to L$ sending $\xi$ to $\eta$.
\end{lemma}

\begin{proof}
Let $U \subset X$ be an affine open mapping into an
affine open of $S$. Then there is an isomorphism
$K|_U \to L|_U$ by Lemma \ref{lemma-relative-dualizing-complex-algebra} and
Dualizing Complexes, Lemma
\ref{dualizing-lemma-uniqueness-relative-dualizing}.
The reader can reuse the argument of that lemma
in the schemes case to obtain a proof in this case.
We will instead use a glueing argument.

\medskip\noindent
Suppose we have an isomorphism $\alpha : K \to L$.
Then $\alpha(\xi) = u \eta$ for some invertible section
$u \in H^0(W, \Delta_*\mathcal{O}_X) = H^0(X, \mathcal{O}_X)$.
(Because both $\eta$ and $\alpha(\xi)$ are generators
of an invertible $\Delta_*\mathcal{O}_X$-module by assumption.)
Hence after replacing $\alpha$ by $u^{-1}\alpha$
we see that $\alpha(\xi) = \eta$.
Since the automorphism group of
$K$ is $H^0(X, \mathcal{O}_X^*)$ by Lemma \ref{lemma-relative-dualizing-RHom}
there is at most one such $\alpha$.

\medskip\noindent
Let $\mathcal{B}$ be the collection of affine opens of
$X$ which map into an affine open of $S$. For each $U \in \mathcal{B}$
we have a unique isomorphism $\alpha_U : K|_U \to L|_U$
mapping $\xi$ to $\eta$ by the discussion in the previous
two paragraphs.
Observe that $\text{Ext}^i(K|_U, K|_U) = 0$ for $i < 0$
and any open $U$ of $X$ by Lemma \ref{lemma-relative-dualizing-RHom}.
By Cohomology, Lemma \ref{cohomology-lemma-vanishing-and-glueing}
applied to $\text{id} : X \to X$ we get a unique morphism
$\alpha : K \to L$ agreeing
with $\alpha_U$ for all $U \in \mathcal{B}$.
Then $\alpha$ sends $\xi$ to $\eta$ as this is true locally.
\end{proof}

\end{document}
